% Generated by explore_fda/make_paper_outputs.R -- regenerate rather than edit
\begin{table}[!htbp]
\centering
\begin{threeparttable}
\caption{Regions of each phase where force was associated with jump height, age and body mass}
\label{tab:fosr}
\footnotesize
\begin{tabular*}{\linewidth}{@{\extracolsep{\fill}}>{\raggedright\arraybackslash}p{4.3cm} >{\raggedright\arraybackslash}p{4.1cm} S[table-format=2.0] S[table-format=-1.3] l@{}}
\toprule
Predictor & {Region (\% of phase)$^{a}$} & {Peak (\%)} & {Effect} & {95\% CI} \\
\midrule
\multicolumn{5}{@{}l}{\textit{Panel A. Take-off, effect in BW per SD of the predictor}} \\
\quad Jump height & {14--30, 66--68, 72--84, 88--98} & 95 & 0.25 & {[0.22, 0.29]} \\
\quad Age & {98} & 98 & 0.09 & {[0.06, 0.12]} \\
\quad Body mass & {2--6, 98--100} & 98 & -0.06 & {[$-$0.09, $-$0.03]} \\
\addlinespace
\multicolumn{5}{@{}l}{\textit{Panel B. Landing, effect in BW per SD of the predictor}} \\
\quad Jump height & {4--28} & 14 & 0.36 & {[0.21, 0.52]} \\
\quad Age & {--} & {--} & {--} & {--} \\
\quad Body mass & {76--82} & 78 & -0.18 & {[$-$0.30, $-$0.06]} \\
\addlinespace
\multicolumn{5}{@{}l}{\textit{Panel C. Jump height split within and between athletes, effect in BW per cm$^{b}$}} \\
\quad Take-off, between athletes & {18--29, 67--99} & 95 & 0.031 & {[0.026, 0.035]} \\
\quad Take-off, within athlete & {3--29, 46--63, 92--99} & 58 & 0.029 & {[0.017, 0.042]} \\
\quad Landing, between athletes & {9--15, 18--25} & 24 & 0.039 & {[0.016, 0.061]} \\
\quad Landing, within athlete & {2--16, 23} & 13 & 0.115 & {[0.074, 0.164]} \\
\bottomrule
\end{tabular*}
\begin{tablenotes}[flushleft]\footnotesize
\item \textit{Note.} Function-on-scalar regression of body-weight-normalised force on jump height, age and body mass, each standardised (1 \textit{SD} = 8.4\,cm, 5.7 years and 15.8\,kg). Effect is the coefficient at the peak of the region, with a pointwise 95\% confidence interval (CI) from 2000 bootstrap resamples of athletes. BW = multiples of body weight; dash = no region.
\item $^{a}$Panels A and B: nodes where the simultaneous 95\% bands of both the athlete-bootstrap least-squares fit and the functional mixed model (athlete and session random curves) exclude zero. $^{b}$Athlete-mean jump height and the trial's deviation from it entered together; region from the athlete-bootstrap simultaneous band.
\end{tablenotes}
\end{threeparttable}
\end{table}
